\subsection{Sections over the Laurent-series field}

We use the arithmetic identification and the basis $(x,y,z)=(b,ab,c)$
of F.9. Put $G=G_{\mathbb Q((t))}$ and
$G_0=G_{\mathbb Q}$. The coefficient splitting fixes all compatible
roots of $t$. A chosen tame generator is denoted by $\tau$;
only its pro-$\ell$ component enters the following formulas.
All cohomology groups and cocycles below are continuous.

\begin{statement}{F.10}{Proved}{An integral tame equation, including the prime two}
Write the class-two free pro-$\ell$ group $P_2$ in Hall coordinates
\[
 x^{u_1}y^{u_2}z^{u_3}
 [x,y]^{w_{12}}[x,z]^{w_{13}}[y,z]^{w_{23}}=(u,w).
\]
These give a homeomorphism of profinite sets
$P_2\simeq\mathbb Z_\ell^3\times\mathbb Z_\ell^3$ with law
\[
 (u,w)(v,h)=(u+v,w+h+F(u,v)),\qquad
 F(u,v)=(-u_2v_1,-u_3v_1,-u_3v_2).
\]
Put $L(u)=(-u_1,0,0)$ and
$B(u,v)=F(u,v)-F(v,u)$.
The arithmetic actions are
\[
 \tau^a(u,w)=(u,w+aL(u)),\qquad
 g(u,w)=(\chi(g)u,\chi(g)^2w)\quad(g\in G_0).
\]
For a section with geometric value $(t,w)$ at $\tau$ and
$(r_g,z_g)$ at $g\in G_0$, the tame relation is exactly
\[
 (\chi(g)-1)w+B(r_g,t)-L(r_g)
 -\frac{\chi(g)-1}{2}\{F(t,t)+L(t)\}=0.                 \tag{F.10.1}
\]
The expression $(\chi(g)-1)/2$ is integral also for $\ell=2$,
because the $2$-adic cyclotomic character is odd.
No Baker--Campbell--Hausdorff denominator is used in this statement.
\end{statement}
\begin{proof}
Moving $x^{v_1}$ and $y^{v_2}$ left through the preceding powers
of $y,z$ gives the displayed $F$ with the convention
$[x,y]=xyx^{-1}y^{-1}$. This proves the group law and,
by uniqueness of Hall coordinates at each finite quotient,
its topological interpretation. Conjugation by $y$ sends
$x$ to $x[x,y]^{-1}$ and fixes $y,z$, proving the inertia action.
The constant action is F.9.

Let $h(t)=F(t,t)+L(t)$. Integer powers of the lifted inertia element
$((t,w),\tau)$ have geometric coordinates
\[
 (nt,nw+\tbinom n2h(t)).
\]
Induction follows from the group law: the next multiplication adds
$nF(t,t)$ and $nL(t)$. The integer-valued polynomial
$\binom n2$ extends continuously to $\mathbb Z_\ell$, giving the same
formula for every $\ell$-adic power. The parts of tame inertia at
other primes act trivially on $P_2$ and have trivial geometric
cocycles, since a continuous homomorphism from a pro-prime-to-$\ell$
group to a pro-$\ell$ group is trivial.

Write $q=\chi(g)$. The equality
$s(g)s(\tau)=s(\tau)^q s(g)$ has equal abelian coordinates
$r_g+qt$ on both sides. Its central defect is
\[
 q(q-1)w+q\{B(r_g,t)-L(r_g)\}
             -\tbinom q2h(t).
\]
Division by the unit $q$ proves (F.10.1).
The program \texttt{f\_tame\_relation.py} checks this identity
symbolically over $\mathbb Q$ and by exact integer arithmetic
in $30\,000$ trials at moduli $8,9,25,27,32,49$.
The symbolic identity and the integral power formula, rather
than those trials, prove the assertion at every coefficient level.
\end{proof}

\begin{statement}{F.11}{Proved}{The continuous obstruction group and all abelian fibers}
There is a split isomorphism of $\mathbb Z_\ell$-modules
\[
 H^2_{\mathrm{cts}}(G,\mathbb Z_\ell(2)^3)
 \simeq
 H^2_{\mathrm{cts}}(G_0,\mathbb Z_\ell(2))^3
 \oplus H^1_{\mathrm{cts}}(G_0,\mathbb Z_\ell(1))^3.       \tag{F.11.1}
\]
The first summand is inflation. A splitting of the second is
cup product with the Kummer class of the uniformizer $t$, with
the orientation chosen so that the residue at $\tau$ is the identity.
The first summand carries a nonzero first-derived-limit term for
every $\ell$. The second carries none.

Let
\[
 \mathscr K_\ell=H^1(G_0,\mathbb Z_\ell(1))
 =
 \begin{cases}
 \mathcal R_{\mathbb P}(\mathbb Z_\ell),&\ell\ne2,\\
 (\mathbb Z/2)\kappa(-1)\oplus
       \mathcal R_{\mathbb P}(\mathbb Z_2),&\ell=2.
 \end{cases}
\]
Let $\epsilon_\ell$ be zero for odd $\ell$ and $\kappa(-1)$ for
$\ell=2$. The abelian section parameters are
$(r,t)\in\mathscr K_\ell^3\times\mathbb Z_\ell^3$.
They lift to $P_2$ exactly when the following six classes vanish:
\[
 r_1\cup r_2,\quad r_1\cup r_3,\quad r_2\cup r_3
       \quad\text{in }H^2_{\mathrm{cts}}(G_0,\mathbb Z_\ell(2)), \tag{F.11.2}
\]
and
\[
 \begin{split}
 &(t_2+1)r_1-t_1r_2+\epsilon_\ell\,t_1(t_2+1),\\
 &t_3r_1-t_1r_3+\epsilon_\ell\,t_1t_3,\\
 &t_3r_2-t_2r_3+\epsilon_\ell\,t_2t_3
 \end{split}
 \qquad\text{in }\mathscr K_\ell.                         \tag{F.11.3}
\]
Every nonempty fiber of geometric conjugacy classes is a torsor under
$H^1_{\mathrm{cts}}(G_0,\mathbb Z_\ell(2))^3$.
In (F.11.2), vanishing means continuous cohomology, computed by
F.1--F.3, including the finite-$S$ term and restricted convergence.
\end{statement}
\begin{proof}
For $M=\mathbb Z_\ell(2)^3$, inertia acts trivially on the coefficient
module. Its cohomology has only degrees zero and one; the latter
is $M(-1)$. The split group extension and the uniformizer Kummer
class give, at every finite coefficient level, inflation and a
right inverse to the residue map by cup product. Its restriction
to inertia is the identity $\mathbb Z_\ell(1)\to\mathbb Z_\ell(1)$.
Thus the two-row Hochschild--Serre complex splits in the indicated
degrees. The same constructions commute with all coefficient
reductions. Derived inverse limits give (F.11.1), so the splitting
also retains the Milnor terms.

For twist two over $\mathbb Q$, the residue differential away from
$\ell$ is multiplication by $q-1$. Its nonzero Smith valuations
are unbounded, by primes $q\equiv1\pmod{\ell^n}$; F.3 gives a
nonzero derived limit. For twist one the residue differential is
zero, so F.3 gives vanishing of its $\varprojlim^1 H^1$ term.
The Milnor term attached directly to the $H^1$ summand in
(F.11.1) is $\varprojlim^1 H^0(G_0,\mathbb Z/\ell^n(1))$;
it vanishes because all these degree-zero groups are finite.
The Kummer description of
$\mathscr K_\ell$ follows directly from
$\mathbb Q^\times=\{\pm1\}\times\bigoplus_p p^{\mathbb Z}$,
together with the finite $H^0$ coefficient comparison.

An abelian cocycle is specified by a cocycle
$r:G_0\to\mathbb Z_\ell(1)^3$ and a vector $t$ on the chosen
inertia generator. The tame relation on this abelian group is
automatic, since both sides scale $t$ by $\chi(g)$.
The central cocycle equation on $G_0$ is
\[
 z_{gh}=z_g+\chi(g)^2z_h+F(r_g,\chi(g)r_h).
\]
Its three cohomology classes are the three cup products in
(F.11.2): the negative reverse cup represents the forward cup,
including at two by the cochain homotopy for graded commutativity.
Vanishing gives continuous choices of $z$.

The remaining equation (F.10.1) is a coboundary equation with
coefficients $\mathbb Z_\ell(1)^3$ for the single vector $w$.
For odd $\ell$, $(\chi-1)/2$ is the coboundary of $1/2$.
For $\ell=2$, the compatible roots of $-1$ identify this cocycle
with $\kappa(-1)$: on a root of unity of order $2^{n+1}$,
its Kummer exponent is $(\chi-1)/2$ modulo $2^n$.
Expanding $B(r,t)-L(r)$ and $F(t,t)+L(t)$ now gives precisely
(F.11.3). Vanishing of these classes supplies an integral $w$.
The cocycle on $G_0$, the lifted inertia generator, and their
verified tame relation define a continuous homomorphism on
$\widehat{\mathbb Z}(1)\rtimes G_0$; the power formula in F.10
proves continuity on inertia. This proves sufficiency as well
as necessity.

For fixed abelian parameters, central lifts differ by
$H^1(G,M)$. Hochschild--Serre gives
$H^1(G,M)=H^1(G_0,M)$, because
$H^0(G_0,M(-1))=0$. There is no further identification by an
abelian geometric stabilizer: $H^0(G,\mathbb Z_\ell(1)^3)=0$.
The central-extension exact sequence therefore gives the asserted
torsor, with the chosen section origin as its sole ambiguity.
\end{proof}

\begin{statement}{F.12}{Proved}{The fully evaluated three-adic Laurent-series test}
For $\ell=3$, write each $r_i=\sum_p r_{i,p}\kappa_p$ as a restricted
family. At $q\ne3$, choose compatible finite-field Kummer logarithms
$\rho_{q,p}$ of the units $p$, for $p\ne q$, modulo $q-1$, and
use integral representatives. Put
\[
 R_{i,q}=\sum_{p\ne q}r_{i,p}\rho_{q,p},\qquad
 c_{ij,q}=-\frac{R_{i,q}r_{j,q}-R_{j,q}r_{i,q}}{q-1}.
\]
The exact conditions (F.11.2) are that each of the three families
$(c_{ij,q})_{q\ne3}$ is integral and restricted.
Condition (F.11.3) is the displayed system of three linear
equations in the restricted Kummer families, with
$\epsilon_3=0$. Every nonempty fiber has $27$ section classes.
\end{statement}
\begin{proof}
The cochain localization calculation of LM7 E.32l--E.32n applies
with commutator coefficient one instead of nine. The comparison
is explicit: the class-two law in logarithmic coordinates has
central term $\tfrac12(u\wedge v)$, so its central tame defect is
$q\{(q-1)c+R_{i,q}r_{j,q}-R_{j,q}r_{i,q}\}$.
Replacing the actual Frobenius Kummer values by the chosen finite-field
representatives adds $(q-1)$ times an integral restricted family:
each difference contains the ramification coefficient
$r_{i,q}$ or $r_{j,q}$. This proves the same residue representative
in the continuous quotient, including its convergence requirement.
For this coefficient, the finite-$S$ $H^2$ is zero by the
independently checked calculation in F.4.

Finally $H^1_{\mathrm{cts}}(G_{\mathbb Q},\mathbb Z_3(2))
\simeq\mathbb Z/3$. One direct verification uses
$H^0(G_S,\mathbb Z/3^n(2))$ of order three, the vanishing of
$H^2(G_S,\mathbb Z/3^n(2))$, and the real Euler characteristic.
They give $|H^1(G_S,\mathbb Z/3^n(2))|=3$ for all $n$.
The nonzero compatible class is
$g\mapsto(\chi_3(g)^2-1)/3$; its class has exact order three,
since $\chi_3^2-1$ has minimum valuation one.
Integral residues in degree one vanish because multiplication
by $q-1$ is injective on $\mathbb Z_3$.
Thus the global and finite-$S$ integral $H^1$ agree and have
the asserted order. The torsor in F.11 has order $3^3=27$.
\end{proof}
